\section{Discussion}
\label{sec:discussion}

The predictable geometry established through Sparse Concept Anchoring opens new paths for model control. Traditional interpretability workflows identify concepts post-hoc, then design interventions around discovered structure. By anchoring concepts during training, we invert this relationship: structure enables intervention by design.

\paragraph{Variability and Fallback Control}
The two intervention modes serve complementary roles (\cref{sec:intervention_tradeoffs}), but weight ablation selectivity varies more (\cref{fig:ablation-boxplots}) because renormalization redistributes removed components. Both would benefit from explicit fallback control: suppression relies on decoder biases (limiting generality), while ablation's redistribution is unreliable. Optimal ablation~\citep{li2024optimalablationinterpretability} could address this by replacing removed components with optimized constants.

\subsection{Limitations and Future Work}

\paragraph{Permanence}
While ablation permanently removes information, it may leave structural traces aiding adversarial recovery. SCA's clean geometry might paradoxically increase this risk: examining encoder activations could reveal smaller pre-normalization values for certain inputs, exposing removed capabilities. Whether this enables practical attacks remains open.

\paragraph{Extension to Other Domains}
Our experiments benefit from RGB data, where concept similarity and labels are
unambiguous. Linguistic domains introduce higher-dimensional manifolds,
context-dependent similarity, and subjective labeling---the fraction of labeled
examples required for anchoring may increase with manifold dimensionality and
curvature. Applying SCA to transformers requires consideration of where to add
geometric constraints---for example, at which layers of the residual
stream---and whether anchoring can resist bypass via residual connections. Our
initial approach might gather concept labels for passages using source metadata
or automated labeling, then apply weak regularization to all tokens in labeled
passages. Given SCA's robustness to label noise this may prove sufficient.

Longer-term validation should target GPT-2--scale transformers and progressively
more abstract concepts. As the number of anchored concepts $K$ grows, the
cumulative effect of competing attraction and repulsion terms may create a
difficult optimization landscape; characterizing both data-requirement scaling
and the point at which anchoring degrades task performance are important open
questions. Multi-dimensional concepts can anchor to subspaces (as with
\concept{vibrant} colors), which may mitigate $K$-scaling by reducing the
effective number of competing terms, while fuzzy concepts might benefit from
weak subspace anchoring that draws related samples to known regions to reduce
the post-hoc search space.
